\begin{claim}{D.2.59}{Proved}{Integral lifts of all seven primary classes}
Put $L=Y\langle7\rangle$ and $E=\Omega^2L$.  With mod-two
coefficients, $\dim H_{13}E=47$, $\dim H_{15}L=33$, and the
subspace
\[
 \mathcal P_{15}=\bigcap_{i>0}\ker\bigl(
      \operatorname{Sq}_*^i:H_{15}L\longrightarrow H_{15-i}L\bigr)
 \tag{D.2.59a}
\]
has dimension seven.  Every element of $\mathcal P_{15}$ is the
reduction of an integral $2$-local homology class.

There is an explicit basis $K_1,\ldots,K_7$ specified below for which
\[
 \begin{gathered}
 \mathcal U=\langle K_1+K_2,K_3,K_4,K_5,K_6,K_7\rangle
   =\mathcal P_{15}\cap\ker\theta,\\
 \langle\theta,K_1\rangle=\langle\theta,K_2\rangle=1,
 \qquad \langle\theta,K_i\rangle=0\quad(3\leq i\leq7).
 \end{gathered}
 \tag{D.2.59b}
\]
Every class in $\mathcal U$ is the reduction of an integral double
homology suspension from $E$.  Neither $K_1$ nor $K_2$ is such an
integral double suspension.  This is an integral-source statement;
it does not, by itself, supply stable spherical representatives
or compute the incoming Adams differential of D.2.58.
\end{claim}
\begin{proof}
All unnamed integral groups in this proof are localized at $2$.
We retain the labels $a,b,c,d,e,h$ dual respectively to
$u_B,v_B,u_B^2,w_B,u_Bv_B,u_Bw_B$, and the classes
$u,D,E_6,S,R,J_1,J_2,P_r,B_1,\ldots,B_7$ of D.2.22 and
D.2.56.  Thus their respective label degrees are
$2,3,4,5,5,7$; the $J_i$ have degree eight and the $B_i,P_r$
have degree ten.

\textbf{The next edge and the double-bar quotient.}
In the Hopf-algebra argument of D.2.56, the new primitive directions
in degree fourteen of $H^*K(\FF_2,4)$ are exactly
\[
 (\operatorname{Sq}^3x_4)^2,\quad
 (\operatorname{Sq}^2\operatorname{Sq}^1x_4)^2,\quad
 \operatorname{Sq}^6\operatorname{Sq}^3\operatorname{Sq}^1x_4.
 \tag{D.2.59c}
\]
Here the first two are the squares of the two degree-seven
generators, and the last is the new Serre polynomial generator.
Let $r:T\to F$ be the actual projection from the space splitting
$T\simeq B\times F$.  On the three actual homology classes
\[
 r_*([v_B^2][u_B^4]),\qquad
 r_*([u_B^2][w_B^2]),\qquad r_*([w_B][t_B]),
\]
the pullbacks of the three classes in \textup{(D.2.59c)} have
pairing matrix
\[
 \begin{pmatrix}1&1&0\\0&1&1\\0&0&1\end{pmatrix}.
 \tag{D.2.59d}
\]
For clarity, this calculation uses $r_*(bc)=P$, $r_*(ad)=Q$,
and $r_*(g)=r_*(h)=r_*(ae)=r_*(a^2b)=0$ in degree seven.
The first two columns are the diagonal pairings with
$(V^*+P^*)^2$ and $(P^*+Q^*)^2$.  In the last column, the reverse
chain $\operatorname{Sq}_*^1\operatorname{Sq}_*^3
\operatorname{Sq}_*^6$ takes the three representatives respectively
through $c^2\mapsto0$, $ax\mapsto ab\mapsto a^2$, and
$bd\mapsto ab\mapsto a^2$.  Thus no multiplicativity of $r$ is
assumed.  This proves injectivity of $p^*$ through degree fourteen.

The first possible kernel degree is now fifteen.  The bounded
free-resolution argument of D.2.56 therefore gives
$t\geq15+4(s-1)$ for positive Tor in homological degree $s$.
Its total degree is at least $14+3(s-1)$, so the actual
Eilenberg--Moore edge remains the quotient edge through degree
thirteen.  Its Hilbert coefficient gives $\dim H_{13}E=47$.

The two ordinary bar filtrations give
\[
 \begin{gathered}
 H_{15}L\cong H_{13}E/\mathcal J,\\
 \begin{split}
 \mathcal J=\langle{}&uH_8E,DS,DR,E_6S,E_6R,\\
       &\{u,S\},\{u,R\},\{D,E_6\},\xi D,\xi E_6\rangle .
 \end{split}
 \end{gathered}
 \tag{D.2.59e}
\]
Here $uH_8E$ contributes five dimensions.  The first nine product
directions are independent by their reduced diagonals in
bidegrees $(5,8)$ and $(6,7)$.  In the first bar, degree fourteen
has their $38$-dimensional quotient and the three additional
cycles $u|S+S|u$, $u|R+R|u$, $D|E_6+E_6|D$.
In the second bar, the four possible length-two terms in total
degree fifteen have independent ordered-product differentials,
as in D.2.56.  The eight products removed from $H_{14}\Omega L$
have three independent leading terms, the three cycles just listed.
Their five remaining terms are exactly the last five directions
of \textup{(D.2.59e)}: the reversed-product differences give the
three Browder brackets, and the two squares give the top
operations \cite[III, Theorem 1.4]{I_CohenLadaMay1976}.
The brackets have distinct ordinary-Lie leading label multisets
$(a,a,g)/(a,a,h)$, $(a,a,a,c)$, and $(a,b,f)$; the two top
operations have distinct outer restricted terms $\xi\{a,b\}$
and $\xi f$.  They are independent modulo the nine products.
Thus $\dim\mathcal J=14$ and $\dim H_{15}L=33$.

\textbf{The actual normalized Browder map.}
We use a reduced ordinary map, including its corrections on
nonprimitive inputs.  A loop-group model of $T$ and the fundamental
circle of its binary little-cubes operation give a loop of maps
$T\times T\to T$ based at multiplication.  Multiply pointwise
by the inverse of that multiplication.  On each axis the resulting
loop comes from unary little-cubes operations, a contractible
space.  Contract these loops and use the homotopy extension
property of the wedge inclusion.  The result is constant on the
axes and at the circle basepoint, and hence gives an actual map
\[
 \mathcal N:\Sigma(T\wedge T)\longrightarrow T.
 \tag{D.2.59f}
\]
Equivalently, it is the double adjoint of the generalized Whitehead
product of the two evaluation maps $\Sigma^2T\to Y$.
The defining circle action is the Browder operation.  The diagonal
and inverse give the precise mod-two formula
\[
 \mathcal N_*(\sigma(x\otimes y))
   =\sum\{x',y'\}\,S_T(x''y''),
 \quad
 \Delta x=\sum x'\otimes x'',\quad
 \Delta y=\sum y'\otimes y''.
 \tag{D.2.59g}
\]
The inverse factor has no circle dependence, so this formula
follows directly by placing the circle fundamental class in the
action factor.  $S_T$ is the antipode of $H_*T$, not that of
$H_*B$.  For example, if $y$ is primitive then
$\mathcal N_*(\sigma(c\otimes y))=\{c,y\}+a\{a,y\}$.

Restrict \textup{(D.2.59f)} along $B\wedge E$ or $B\wedge B$.
The suspended sources are respectively seven- and four-connected;
their maps to $T$ lift uniquely to the four-connected cover $E$.
Double suspension and evaluation consequently give actual maps
\[
 \Sigma^3(B\wedge E)\longrightarrow L,
 \qquad \Sigma^3(B\wedge B)\longrightarrow L.
 \tag{D.2.59h}
\]
Write $n(x,y)$ for the induced degree-fifteen class when
$|x|+|y|=12$.

\textbf{The seven primary classes.}
An explicit choice of basis for \textup{(D.2.59a)} is
\[
 \begin{aligned}
 K_1&=n(a,B_7),&K_2&=n(a,B_3),\\
 K_3&=n(a,B_2+B_6),&K_4&=n(a,B_2)+n(c,J_2),\\
 K_5&=n(a,P_r),\\
 K_6&=n(a,B_2+B_4)+n(c,J_1)+n(b,K_{\rm sum}),\\
 K_7&=n(e,h),&K_{\rm sum}&=K_2^{\rm old}+K_3^{\rm old}+K_4^{\rm old}.
 \end{aligned}
 \tag{D.2.59i}
\]
The superscript distinguishes the earlier degree-nine cover
classes from the present degree-fifteen basis.
We specify the finite rank verification to make completeness
checkable.  In the ordered quotient basis of
\textup{(D.2.59e)}, the seven vectors have supports
\[
 \begin{array}{c|l}
 K_1&05\\K_2&08\\K_3&06,10\\K_4&06,12\\K_5&15\\
 K_6&06,07,11,13,18,19\\K_7&30.
 \end{array}
 \tag{D.2.59j}
\]
These are the indices $\texttt{L15\_01},\ldots,\texttt{L15\_33}$
of the complete $33$-vector quotient basis in the certificate.
Applying Cartan, the Serre label operations, and the CLM
restricted-operation and Browder formulas gives rank $26$ for
the combined map $(\operatorname{Sq}_*^1,
\operatorname{Sq}_*^2,\operatorname{Sq}_*^4)$.
Its kernel consists exactly of \textup{(D.2.59j)}.
$\operatorname{Sq}_*^8$ is zero by instability on degree-seven
cohomology, and the squares of power-of-two degree generate the
Steenrod algebra.  This proves the asserted primary dimension.
The complete basis, operation matrices, and exact row reductions
are in \texttt{I\_E2\_primary\_15.py},
\texttt{I\_L15\_primary.json}, and
\texttt{I\_L15\_primary\_cycles.json}; the normalized image
calculation is \texttt{I\_integral\_browder\_sources.py}.
These use the actual edge and quotient proved above.  They also
check membership in $H_*E$ before taking the bar quotient,
and all Cartan, coproduct, and Adem compatibilities in this range.

\textbf{Prescribed integral tensors and Tor cycles.}
The classes $a,c,d+e,h$ in $B$ are reductions of integral cycles;
the first two have orders $2,4$, and the last two have order two.
For example $d+e$ and $h$ are reductions of the integral
Bocksteins of $[v_B^2]$ and $[v_Bw_B]$.
The classes $J_1,J_2$ in $E$ also lift integrally:
$H_7(E;\mathbb Z_{(2)})=\mathbb Z/2$, so reduction is injective
on the possible integral connecting image, and their first
Bocksteins vanish.  In degree ten the needed integral reductions
are $P_r,B_2,B_1+B_3,B_5+B_6$.  Indeed
\[
 H_9(E;\mathbb Z_{(2)})=
       \mathbb Z/8\{N'\}\oplus(\mathbb Z/2)^3.
 \tag{D.2.59k}
\]
A first-Bockstein cycle can therefore have integral connecting
image only $0$ or $4N'$.  The latter is detected by the third
Bockstein.  The complete D.2.22 table has
$\beta_3B_5=\beta_3B_6=N$ and zero on
$P_r,B_2,B_1+B_3$, proving the assertion for the four combinations.
Thus no general inference from a first-Bockstein cycle to an
integral lift is used.

Choose integral chains with the following boundaries and the
indicated mod-two reductions:
\[
 \begin{gathered}
 d\widetilde b=2\widetilde a,
 \quad d\widetilde d=4\widetilde c,
 \quad d\widetilde l=2\widetilde h,\\
 d\widetilde B_5=8\widetilde N,
 \qquad d\widetilde B_4=2\widetilde K_{\rm sum}.
 \end{gathered}
 \tag{D.2.59l}
\]
Here $l=[v_Bw_B]$.  These are the established bottom Moore,
second-Bockstein, and D.2.22 Bockstein representatives, with
orientations chosen to give the displayed signs.  The following
literal tensor chains are cycles:
\[
 \begin{aligned}
 \tau(a,N)&=4\widetilde b\otimes\widetilde N
                 -\widetilde a\otimes\widetilde B_5,
       &\rho_2\tau(a,N)&=a\otimes B_5,\\
 \tau(a,K_{\rm sum})&=\widetilde b\otimes\widetilde K_{\rm sum}
                 -\widetilde a\otimes\widetilde B_4,
       &\rho_2\tau(a,K_{\rm sum})&=b\otimes K_{\rm sum}+a\otimes B_4,\\
 \tau(c,h)&=\widetilde d\otimes\widetilde h
                 -2\widetilde c\otimes\widetilde l,
       &\rho_2\tau(c,h)&=d\otimes h.
 \end{aligned}
 \tag{D.2.59m}
\]
The boundaries cancel as $8aN-8aN$, $2aK_{\rm sum}-2aK_{\rm sum}$,
and $4ch-4ch$.  This proves the prescribed reductions directly,
without selecting an abstract K\"unneth splitting.

The actual maps \textup{(D.2.59h)} applied to the following
integral tensor/Tor cycles give six independent directions:
\[
 \begin{array}{c|c}
 \text{integral source cycle}&\text{mod-two image in }H_{15}L\\ \hline
 a\otimes P_r&K_5\\
 a\otimes(B_1+B_3+B_2)&K_1+K_2\\
 a\otimes B_2+c\otimes J_2&K_4\\
 a\otimes(B_2+B_5+B_6)+\tau(a,N)&K_3\\
 a\otimes B_2+c\otimes J_1+\tau(a,K_{\rm sum})&K_6\\
 (d+e)\otimes h+\tau(c,h)&K_7.
 \end{array}
 \tag{D.2.59n}
\]
Integral representatives are understood in the first column.
The first five rows use $B\wedge E$ and the last $B\wedge B$.
The fourth reduction is $a\otimes(B_2+B_6)$ and the last is
$e\otimes h$.  Formula \textup{(D.2.59g)} verifies every image;
in particular its correction on $c$ is essential in the third
and fifth rows.  Independence follows at once from
\textup{(D.2.59j)}.

\textbf{The integral double-suspension restriction.}
By D.2.50, $\beta_{\mathbb Z}\theta=b_8^2$.
Cohomology suspension kills products, so $\sigma^2\theta$ has
an integral cohomology lift on $E$.  The space $E$ is rationally
acyclic of finite type, hence $H_{13}(E;\mathbb Z_{(2)})$ is
torsion.  An integral cohomology class pairs trivially with this
group.  Therefore $\theta$ vanishes on every integral double
homology suspension from $E$, in particular on all six rows of
\textup{(D.2.59n)}.

D.2.52 supplies an integral class $\omega$ whose reduction lies
in $\mathcal P_{15}$ and pairs to one with $\theta$.
Its reduction is outside the six-dimensional subspace
$\mathcal U$.  Thus the reductions of integral classes span all
seven dimensions of $\mathcal P_{15}$.  The functional $\theta$
has kernel exactly $\mathcal U$, proving \textup{(D.2.59b)}.
The nonzero evaluations on $K_1,K_2$ also prove that neither is
an integral double suspension, despite mod-two double-suspension
surjectivity.  The argument imposes no stable-sphericity claim.
\end{proof}
